\section{MODULES OVER INTEGRALLY CLOSED NOETHERIAN DOMAINS}

Throughout this paragraph, A will be a commutative integral domain with field of fractions K. Starting with no. 2, A will be further assumed to be Noetherian and integrally closed (and hence a Krull domain (§ 1, no. 3, Corollary to Theorem 2)); then P(A), D(A) and C(A) will respectively denote the set of prime ideals of A of height 1 (§ 1, no. 6), the divisor group of A (§ 1, no. 3) and the divisor class group of A (§ 1, no. 10), these latter being written additively.

The general method of studying finitely generated modules over an integrally closed Noetherian domain A consists of ``localizing'' the modules with respect to all the prime ideals \(p \in \mathrm{P}(\mathrm{A})\) of height 1 in A; as \(A_{p}\) is then a discrete valuation ring (§ 1, no. 6, Theorem 4), the structure of finitely generated \(A_{p}\) -modules is well known (Algebra, Chapter VII, § 4) and therefore gives information about the structure of finitely generated A-modules. In the particular case where A is a Dedekind domain, we can arrive at as complete a theory as when A is a principal ideal domain (no. 10).

